\documentclass[../DoAn.tex]{subfiles}
\begin{document}

\section{Giải pháp multi-tenant RAG chatbot cho tư vấn bán hàng nội thất}
\label{section:5.1}

\subsection{Bài toán đặt ra}
\label{subsection:5.1.1}

Trong các chương trước, đồ án đã xác định bài toán chính là xây dựng một hệ thống chatbot hỗ trợ tư vấn bán hàng nội thất cho nhiều cửa hàng trên cùng một nền tảng. Điểm khó của bài toán không chỉ nằm ở việc xây dựng một chatbot có khả năng trả lời bằng ngôn ngữ tự nhiên, mà còn ở việc tổ chức hệ thống sao cho mỗi cửa hàng có dữ liệu, chatbot, hội thoại và yêu cầu mua hàng riêng. Nếu hệ thống chỉ có một chatbot dùng chung cho tất cả cửa hàng, chatbot có thể trả lời lẫn dữ liệu giữa các cửa hàng, gây sai lệch trong quá trình tư vấn. Ngược lại, nếu xây dựng một hệ thống riêng cho từng cửa hàng, việc triển khai, vận hành và bảo trì sẽ trở nên phức tạp khi số lượng cửa hàng tăng lên.

Bên cạnh đó, dữ liệu sản phẩm nội thất có tính thay đổi theo từng cửa hàng. Mỗi cửa hàng có danh mục sản phẩm, cách mô tả, mức giá và chính sách tư vấn riêng. Vì vậy, chatbot không thể chỉ dựa vào kiến thức chung của mô hình ngôn ngữ để trả lời, mà cần có khả năng truy xuất thông tin từ dữ liệu của đúng cửa hàng đang được tư vấn. Đây là lý do hệ thống cần kết hợp mô hình multi-tenant với hướng tiếp cận Retrieval-Augmented Generation. Multi-tenant giúp tổ chức nhiều cửa hàng trên cùng một nền tảng nhưng vẫn tách biệt dữ liệu, còn RAG giúp chatbot trả lời dựa trên dữ liệu được truy xuất từ knowledge base của từng cửa hàng.

Một vấn đề khác là quá trình tư vấn bán hàng không kết thúc ở việc chatbot trả lời câu hỏi. Trong thực tế, khi khách hàng quan tâm đến sản phẩm hoặc có ý định mua hàng, cửa hàng cần tiếp nhận thông tin đó để tiếp tục xử lý. Nếu chatbot chỉ dừng lại ở trả lời hội thoại, hệ thống sẽ thiếu liên kết với nghiệp vụ bán hàng phía sau. Vì vậy, đồ án cần thiết kế một luồng xử lý trong đó hội thoại có thể phát sinh lead hoặc yêu cầu mua hàng, giúp cửa hàng theo dõi và xử lý tiếp sau khi khách hàng đã trao đổi với chatbot.

Từ các vấn đề trên, đóng góp chính của phần này là đề xuất và xây dựng một giải pháp kết hợp giữa mô hình multi-tenant, chatbot RAG và luồng nghiệp vụ bán hàng. Giải pháp này hướng đến việc giải quyết đồng thời ba yêu cầu: tư vấn dựa trên dữ liệu riêng của từng cửa hàng, tách biệt dữ liệu giữa các tenant, và chuyển kết quả hội thoại thành thông tin nghiệp vụ để cửa hàng tiếp tục xử lý.

\subsection{Giải pháp đề xuất}
\label{subsection:5.1.2}

Giải pháp được đề xuất là tổ chức hệ thống theo mô hình multi-tenant RAG chatbot. Trong mô hình này, mỗi cửa hàng được biểu diễn như một tenant trong hệ thống. Các thực thể quan trọng như chatbot, nguồn dữ liệu knowledge base, hội thoại, lead và yêu cầu mua hàng đều được gắn với tenant tương ứng. Nhờ đó, hệ thống có thể phục vụ nhiều cửa hàng trên cùng một nền tảng, nhưng vẫn đảm bảo dữ liệu của từng cửa hàng được quản lý riêng biệt.

Ở tầng nghiệp vụ, hệ thống cung cấp các chức năng quản lý tenant, quản lý thành viên cửa hàng, quản lý chatbot, quản lý nguồn dữ liệu knowledge base, quản lý hội thoại, quản lý lead và quản lý yêu cầu mua hàng. Quản trị hệ thống có vai trò tạo và quản lý các tenant, trong khi quản trị cửa hàng chịu trách nhiệm cấu hình chatbot và nguồn dữ liệu của cửa hàng mình. Nhân viên cửa hàng theo dõi hội thoại, lead và yêu cầu mua hàng phát sinh từ quá trình tư vấn. Cách phân chia này giúp hệ thống phản ánh đúng các vai trò trong bài toán thực tế, đồng thời hỗ trợ kiểm soát quyền truy cập theo tenant.

Ở tầng xử lý hội thoại, khi người dùng gửi tin nhắn đến chatbot của một cửa hàng, hệ thống trước hết xác định chatbot và tenant tương ứng. Sau đó, yêu cầu hội thoại được chuyển đến AI service cùng với thông tin ngữ cảnh cần thiết. AI service thực hiện truy xuất dữ liệu liên quan từ knowledge base của tenant đó, rồi sử dụng kết quả truy xuất làm ngữ cảnh để sinh câu trả lời tư vấn. Cách xử lý này giúp chatbot phản hồi dựa trên dữ liệu của đúng cửa hàng, thay vì trả lời chung chung hoặc sử dụng dữ liệu không thuộc phạm vi tư vấn.

Knowledge base được thiết kế như thành phần trung gian giữa dữ liệu cửa hàng và chatbot. Thay vì yêu cầu cửa hàng tự xây dựng hệ thống truy xuất phức tạp, hệ thống cung cấp chức năng quản lý nguồn dữ liệu và rebuild knowledge base. Khi cửa hàng cập nhật nguồn dữ liệu, hệ thống có thể xử lý lại dữ liệu để phục vụ truy xuất trong các phiên tư vấn sau. Nhờ đó, quy trình cập nhật tri thức cho chatbot được đưa về một nhóm thao tác quản trị rõ ràng, phù hợp hơn với người dùng phía cửa hàng.

Một điểm quan trọng trong giải pháp là liên kết giữa hội thoại và nghiệp vụ bán hàng. Trong quá trình tư vấn, nếu người dùng thể hiện ý định mua hàng hoặc cung cấp thông tin liên hệ, hệ thống có thể tạo lead hoặc yêu cầu mua hàng từ hội thoại. Các bản ghi này lưu lại thông tin người dùng, nhu cầu, sản phẩm quan tâm và nội dung hội thoại liên quan. Sau đó, nhân viên cửa hàng có thể truy cập danh sách lead hoặc yêu cầu mua hàng để claim, phản hồi, phân công hoặc cập nhật trạng thái xử lý. Nhờ vậy, chatbot không chỉ là công cụ trả lời tự động, mà còn trở thành điểm đầu vào cho quy trình bán hàng của cửa hàng.

Giải pháp cũng tách riêng backend nghiệp vụ và AI service. Backend chịu trách nhiệm xác thực, phân quyền, tenant isolation, quản lý dữ liệu nghiệp vụ và giao tiếp với frontend hoặc kênh ngoài. AI service tập trung vào xử lý hội thoại, truy xuất knowledge base và sinh câu trả lời. Cách tách này giúp hệ thống dễ bảo trì hơn, vì phần nghiệp vụ và phần AI có thể được phát triển, kiểm thử hoặc thay đổi tương đối độc lập. Khi cần thay đổi mô hình ngôn ngữ, chiến lược truy xuất hoặc cách xây dựng prompt, phần thay đổi chủ yếu nằm ở AI service mà không ảnh hưởng lớn đến các chức năng quản trị nghiệp vụ.

Đối với người dùng cuối, giải pháp cung cấp một trải nghiệm tương đối liền mạch: người dùng có thể trao đổi với chatbot, bổ sung hoặc thay đổi nhu cầu, nhận tư vấn và gửi yêu cầu mua hàng khi cần. Đối với cửa hàng, giải pháp cung cấp các chức năng để vận hành chatbot trên dữ liệu riêng, theo dõi kết quả tư vấn và xử lý các cơ hội bán hàng phát sinh. Đối với quản trị hệ thống, giải pháp hỗ trợ tổ chức nhiều cửa hàng trên cùng một nền tảng, phù hợp với định hướng multi-tenant của đồ án.

\subsection{Kết quả đạt được}
\label{subsection:5.1.3}

Kết quả đạt được của giải pháp là một thiết kế hệ thống có khả năng kết hợp ba thành phần chính: multi-tenant, RAG chatbot và quy trình xử lý lead/yêu cầu mua hàng. Về mặt tổ chức dữ liệu, hệ thống xác định tenant là đơn vị trung tâm để phân tách dữ liệu giữa các cửa hàng. Các chatbot, nguồn dữ liệu knowledge base, hội thoại, lead và yêu cầu mua hàng đều được gắn với tenant tương ứng. Cách tổ chức này giúp hệ thống phục vụ được nhiều cửa hàng mà không cần triển khai riêng một hệ thống cho từng cửa hàng.

Về mặt tư vấn, hệ thống sử dụng hướng tiếp cận RAG để chatbot có thể trả lời dựa trên dữ liệu riêng của cửa hàng. Khi người dùng hỏi về sản phẩm, hệ thống không chỉ gửi câu hỏi trực tiếp cho mô hình ngôn ngữ, mà còn truy xuất thông tin liên quan từ knowledge base để bổ sung ngữ cảnh trả lời. Điều này giúp chatbot có khả năng bám sát dữ liệu cửa hàng hơn so với một chatbot chỉ trả lời dựa trên kiến thức chung.

Về mặt nghiệp vụ, hệ thống không dừng lại ở hội thoại mà còn ghi nhận các thông tin có giá trị cho cửa hàng. Khi người dùng có ý định mua hàng hoặc để lại thông tin liên hệ, hệ thống có thể tạo lead hoặc yêu cầu mua hàng. Nhân viên cửa hàng có thể tiếp nhận các thông tin này, xem lại nội dung hội thoại và cập nhật trạng thái xử lý. Kết quả này giúp liên kết giữa quá trình tư vấn tự động và quy trình bán hàng thủ công của cửa hàng.

Về mặt kiến trúc, việc tách backend nghiệp vụ và AI service giúp hệ thống rõ ràng hơn về trách nhiệm của từng thành phần. Backend tập trung vào quản lý dữ liệu, phân quyền và nghiệp vụ multi-tenant; AI service tập trung vào xử lý hội thoại và truy xuất tri thức. Cách tổ chức này cũng tạo điều kiện để hệ thống có thể mở rộng trong tương lai, ví dụ bổ sung kênh chat ngoài, thay đổi mô hình ngôn ngữ, cải thiện thuật toán truy xuất hoặc mở rộng chức năng quản lý knowledge base.

Nhìn chung, đóng góp của phần này không nằm ở việc xây dựng một nền tảng AI agent tổng quát, mà nằm ở việc thiết kế và hiện thực một giải pháp phù hợp với bài toán hẹp của đồ án: tư vấn bán hàng nội thất cho nhiều cửa hàng trên cùng một nền tảng, trong đó mỗi cửa hàng có dữ liệu và chatbot riêng, còn kết quả hội thoại có thể được chuyển thành lead hoặc yêu cầu mua hàng để cửa hàng tiếp tục xử lý.

% Các phần đóng góp khác sẽ được bổ sung sau nếu cần:
% \section{Giải pháp quản lý knowledge base và rebuild dữ liệu theo tenant}
% \section{Giải pháp chuyển hội thoại tư vấn thành lead/yêu cầu mua hàng}
% \section{Giải pháp tham chiếu giá cơ bản trong phạm vi dữ liệu hệ thống}

\end{document}